
\chapter{Projeto Eletrônico}
\label{chap:projeto-eletronico}

% Apresentar o desenvolvimento eletrônico do DoseCerta,
% relacionando os circuitos e componentes aos requisitos do sistema.

\section{Requisitos Eletrônicos}
\label{sec:requisitos-eletronicos}

% Descrever os requisitos considerados no projeto:
% alimentação, controle, sensores, atuadores e interfaces.
% Relacionar, quando possível, aos requisitos do Capítulo 5.

\section{Arquitetura Eletrônica}
\label{sec:arquitetura-eletronica}

% Apresentar a organização dos módulos e componentes eletrônicos.
% Inserir diagrama de blocos, se disponível.

\section{Unidade de Processamento}
\label{sec:unidade-processamento}

% Identificar a placa ou unidade de processamento utilizada.
% Justificar a escolha com base nas necessidades do projeto.
% Detalhes de programação e execução ficam no Capítulo 9.

\section{Sensores}
\label{sec:sensores}

% Para cada sensor utilizado, descrever:
% função no sistema, princípio de funcionamento,
% interface elétrica e critérios de escolha.
% Inserir tabela ou figura, se ajudar na compreensão.

\section{Atuadores}
\label{sec:atuadores}

% Descrever os atuadores empregados no dispensador,
% sua função, características relevantes e interface de acionamento.
% Explicar a relação com o mecanismo descrito no Capítulo 7.

\section{Circuitos de Alimentação e Acionamento}
\label{sec:alimentacao-acionamento}

% Descrever a alimentação dos módulos e os circuitos de acionamento.
% Informar tensões e correntes somente quando confirmadas
% por documentação, medição ou especificação dos componentes.

\section{Interligações e Comunicação entre Módulos}
\label{sec:interligacoes}

% Apresentar as conexões elétricas entre os componentes.
% Descrever os protocolos ou sinais efetivamente utilizados.
% Incluir esquema elétrico ou diagrama de interconexão, se disponível.

\section{Montagem e Integração Eletrônica}
\label{sec:montagem-integracao-eletronica}

% Descrever a montagem física dos circuitos e sua integração
% ao protótipo. Registrar placas, conectores e adaptações realizadas.

\section{Considerações de Segurança e Confiabilidade}
\label{sec:seguranca-eletronica}

% Discutir aspectos pertinentes ao projeto, como proteção elétrica,
% estabilidade da alimentação, falhas de conexão e comportamento
% seguro em caso de falha, conforme aplicável.

\section{Síntese do Projeto Eletrônico}
\label{sec:sintese-eletronica}

% Resumir as principais escolhas eletrônicas e como elas
% viabilizam as funções do DoseCerta.
